\documentclass[11pt]{article}
\usepackage{amsmath,amsfonts,mathabx}
\usepackage{graphicx}
\usepackage{hyperref}

% PREAMBLE
\newcommand{\RR}{\ensuremath{\mathbb{R}}}
\newcommand{\RRn}[1]{\mathbb{R}^{#1}}
\DeclareMathOperator{\function}{function}
 
% Code listing
\usepackage[listings]{tcolorbox}
\usepackage[normalem]{ulem}
\usepackage{xcolor}


\begin{document}


\begin{center}
{\textbf{\LARGE Chapitre 5: Mod\`eles lin\'eaires g\'en\'eralis\'es}}
\end{center}

\noindent Jusqu'\`a maintenant, on a vu la r\'egression lin\'eaire (Chapitre 2 et 3) et la r\'egression logistique (Chapitre 4). On constate que plusieurs concepts comme l'inf\'erence statistique concernant les coefficients de r\'egression $\beta$, la s\'election de mod\`eles et la d\'etection sont identiques pour les deux situations, m\^eme si chaque mod\`ele peut avoir des particularit\'es qui requi\`erent des proc\'edures sp\'ecifiques comme la courbe ROC pour la r\'egression \sout{lin\'eaire}\textcolor{red}{logistique}. Dans ce chap\sout{ire}\textcolor{red}{itre}, on consid\`ere une pr\'esentation unifi\'ee pour ces mod\`eles. On suppose que la variable r\'eponse $Y_i$ suit une distribution qui appartient \`a la famille exponentielle qui inclut la loi normale, la loi de Bernoulli et la loi binomiale.

\section{Famille exponentielle}

\subsection{D\'efinition et exemples}

Cette famille de lois contient la vaste majorit\'e des lois couramment utilis\'ees en pratique. On dit que la distribution de la variable al\'eatoire $Y$ appartient \textcolor{red}{\`a} la famille exponentielle si la fonction de probabilit\'e de $Y$ s'\'ecrit sous la forme:
$$f_Y(y) = a(y,\phi) \exp\left\{\frac{y\theta-b(\theta)}{\phi}\right\}.$$
\noindent Ici, $f_Y(y)$ est la fonction de densit\'e si $Y$ est continue et la fonction de mas\sout{s}\textcolor{red}{se} $P(Y=y)$ si $Y$ est discr\`ete.

\noindent La distribution est param\'etr\'ee avec deux param\`etres: le param\`etre canonique $\theta$ et le param\`etre de dispersion $\phi$. On montre que $E[Y]=\sout{\mu }b'(\theta)$, o\`u $b'(u)=db(u)/du$ est la d\'eriv\'ee de $b(u)$ par rapport \`a $u$. Par cons\'equent, le param\`etre canonique est directement li\'e \`a l'esp\'erance de $Y$. On montre aussi que $Var(Y)=\phi b''(\theta)$, o\`u $b''(u)=d^2b(u)/du^2$ est la d\'eriv\'ee seconde de $b(u)$ par rapport \`a $u$. On en d\'eduit que $Var[Y]=\phi V(\mu)$, avec $V(u)=b''[{b'}^{-1}(u)]$. Le param\`etre de dispersion $\phi$ mod\'elise donc la variance de $Y$.

\noindent Les distributions de la famille exponentielle poss\`\textcolor{red}{e}dent, en g\'en\'eral, deux param\`etres $\theta$ et $\phi$. Cependant, pour certaines distributions, le param\`etre de dispersion est fix\'e \`a 1 et on a ai\sout{sn}\textcolor{red}{n}si un seul param\`etre.

\noindent Voici quelques exemples de ces distributions:
\begin{enumerate}
\item[$\bullet$] {\bf Loi normale:} Soit $Y \sim {\cal N}(\mu,\sigma^2)$. La fonction de densit\'e de $Y$ s'\'ecrit:
\begin{eqnarray*}
F_Y(y)&=&\frac{1}{\sqrt{2\pi \sigma^2}} \exp\left\{-\frac{(y-\mu)^2}{2\sigma^2}\right\} \\
&=& \frac{1}{\sqrt{2\pi \sigma^2}} \exp\left\{\frac{y\mu-\mu^2/2}{\sigma^2}-\frac{y^2}{2\sigma^2}\right\}
\end{eqnarray*}

Cette fonction de densit\'e est un cas particulier d'une fonction de probabilit\'e d'une loi appartenant \`a la famille exponentielle avec $\theta=\mu$, $\phi=\sigma^2$, $$a(\sout{\theta}\textcolor{red}{y,\phi})=\frac{1}{\sqrt{2\pi \sigma^2}} \exp\left\{-\frac{y^2}{2\sigma^2}\right\}$$ et $b(\theta)=\theta^2/2$. Plus encore, on a $E(Y)=b'(\theta)=\theta=\mu$ et $Var(Y)=\phi b''(\theta) = \sigma^2$ puisque $b''(u)=1$. On en d\'eduit que $V(\mu)=1$.

\item[$\bullet$]{\bf Loi binomiale:} Soit $Y \sim \mbox{Bin}(m,\pi)$. On a alors pour $y=0,1,\cdots,m$:
\begin{eqnarray*}
P(Y=y) &=& \binom{m}{y} \pi^y (1-\pi)^{m-y}\\
&=& \binom{m}{y} \left(\frac{\pi}{1-\pi}\right)^y (1-\pi)^m \\
&=& \binom{m}{y} \exp\left\{y \log\left(\frac{\pi}{1-\pi}\right) +m\log(1-\pi)\right\} \\
\end{eqnarray*}

On reconnait la fonction de probabilit\'e d'une famille exponentielle avec:
\begin{eqnarray*}
\theta &=& \log\left\{\frac{\pi}{1-\pi}\right\} \\
\pi&=& \frac{e^{\theta}}{1+e^{\theta}} \\
\phi&=&1 \\
a(y,\phi)&=& \binom{m}{y} \\
b(\theta)&=&-m\log(1-\pi) = m\log(1+e^{\theta}).
\end{eqnarray*}

On en d\'eduit que $E[Y]=b'(\theta)=m e^{\theta}/(1+e^{\theta})=m\pi$, $Var[Y]=\phi b''(\theta)=me^{\theta}/(1+e^{\theta})^2=m\pi(1-\pi)$ et $V(\mu)=\mu(m-\mu)/m$.

\item[$\bullet$] {\bf Loi de Poisson:} Soit $Y \sim {\cal P}oisson(\lambda)$. On a alors pour $y=0,1,\cdots$, on a $$P(Y=y)=e^{-\lambda} \frac{\lambda^y}{y!}=\frac{1}{y!} \exp\left\{y\log(\lambda)-\lambda\right\}.$$ On reconnait la fonction de probabilit\'e d'une famille exponentielle avec: $\theta=\log(\lambda)$, $\phi=1$, $a(y,\phi)=1/y!$ et $b(\theta)=\lambda = e^{\theta}.$ On en d\'eduit que $E[Y]=b'(\theta)=e^{\theta}=\lambda=\mu$, $Var[Y]=\phi b''(\theta)=e^{\theta}=\lambda=\mu$ et $V(\mu)=\mu$.


\item[$\bullet$] {\bf Loi Gamma:} Soit $Y \sim \mbox{Gamma}(\alpha,\beta)$. La fonction de densit\'e de $Y$ est donn\'ee par 
$$
f_Y(y) = \begin{cases} \frac{y^{\alpha-1}e^{-y/\beta}}{\Gamma(\alpha)\beta^\alpha}  & \mbox{si } y \ge 0 \\
0 & \mbox{si } y<0
\end{cases}
$$
o\`u $\Gamma(\alpha)=\int_0^\infty t^{\alpha-1} e^{-t} dt$. Avec cette param\'etrisation, on a $E[Y]=\alpha\beta$ et $Var[Y]=\alpha \beta^2$. Pour montrer que la loi Gamma appartient \`a la famille exponentielle, on reparam\'etrise la fonction de densit\'e comme suit. On pose $\mu=\alpha\beta$ et $\phi=1/\alpha$ et on obtient
\begin{eqnarray*}
f_Y(y) &=& \left(\frac{y}{\phi\mu}\right)^{1/\phi} \frac{1}{y} \frac{1}{\Gamma(1/\phi)} \exp\left\{-\frac{y}{\phi\mu}\right\}  \\
&=& \left(\frac{y}{\mu}\right)^{1/\phi} \frac{1}{y} \frac{1}{\Gamma(1/\phi)} \exp\left\{-\frac{y}{\phi\mu}-\frac{\log(\mu)}{\phi}\right\}.
\end{eqnarray*}

On reconnait la fonction de probabilit\'e d'une famille exponentielle avec:
\begin{eqnarray*}
\theta &=& -1/\mu \\
a(y,\phi)&=& \left(\frac{y}{\mu}\right)^{1/\phi} \frac{1}{y} \frac{1}{\Gamma(1/\phi)}\\
b(\theta)&=&\log(\mu)=\log(-1/\theta)=-\log(-\theta)
\end{eqnarray*}

On en d\'eduit que $E[Y]=b'(\theta)=-1/\theta=\mu$, $Var[Y]=\phi b''(\theta) = \phi/\theta^2 = \phi \mu^2$ et $V(\mu)=\mu^2$.
\end{enumerate}

\subsection{D\'eviance}

Soit $Y$ une variable a\'eatoire qui suit une distribution appartenant \`a la famille exponentielle. On pose $t(y,\mu)=y\theta-b(\theta)$ et $d(y,\mu)=2\{t(y,y)-t(y,\mu)\}$. On a alors 
\begin{eqnarray*}
\frac{\partial d(y,\mu)}{\partial \theta}&=& - \left\{y - b'(\theta)\right\} \\
\frac{\partial^2 d(y,\mu)}{\partial \theta^2}&=&   b''(\theta) = V(\mu) >0 \\
\end{eqnarray*}
\noindent On en d\'eduit que $\partial d(y,\mu)/\partial \theta$ est \'egale \`a 0 lorsque $\mu=y$. Par cons\'equent, $d(y,\mu)$ est toujours positive et ne s'annule que lorsque $\mu=y$. 

\noindent Pour la loi normale ${\cal N}(\mu,\sigma^2)$, on a $t(y,\mu)=y\mu-\mu^2/2$ et $d(y,\mu)=(y-\mu)^2$. 

\noindent \`A la table 5.1 du livre {\it Generalized linear models with examples in R}, on trouve les expressions de $d(y,\mu)$ pour plusieurs distributions usuelles.

\noindent On montre que lorsque $Y$ suit une distribution appartenant \`a la famille exponentielle, $d(Y,\mu)/\phi$ suit approximativement une distribution de Khi-deux \`a 1 degr\'e de libert\'e $\chi^2_1$. Supposons maintenant qu'on a un ensemble de variables al\'eatoires ind\'ependantes $\{Y_1,\cdots,Y_n\}$ qui suivent une distribution appartenant \`a la famille exponentielle. On suppose que pour $i=1,\cdots,n$, les param\`etres de la distributions de $Y_i$ sont $(\mu_i,\phi)$. Ainsi, chaque variable al\'eatoire poss\`ede sont propre param\`etre canonique mais le param\`etre de dispersion $\phi$ est commun \`a l'ensemble des variables. La d\'eviance et la d\'eviance standardi\'ee de $\{Y_1,\cdots,Y_n\}$ sont d\'efinies respectivement par 
$D(y,\mu) = \sum_{i=1}^n d(y_i,\mu_i)$ et $D^*(y,\mu)=D(y,\mu)/\phi$. On en d\'eduit que $D^*(y,\mu)$ suit approximativement une distribution de Khi-deux \`a $n$ degr\'es de libert\'e $\chi^2_n$.

\section{Mod\`eles lin\'eaires g\'en\'eralis\'es}

\subsection{Donn\'ees et mod\`ele}
Consid\'erons les donn\'ees $\Delta=\{(Y_i,X_i), i=1,\cdots,n\}$. Pour $i=1,\cdots,n$, $Y_i$ est la variable r\'eponse et $X_i=(X_{i1},\cdots,X_{ip})^\top$ est un ensemble de $p$ variables explicatives pour l'observation $i$. Un mod\`ele lin\'eaire g\'en\'earlis\'e est compos\'e de 3 composantes:
\begin{enumerate}
\item[$\bullet$] {\bf Composante al\'eatoire:} On sp\'ecifie une distribution pour $Y_i$. Explicitement, on suppose que $Y_i$ suit une distribution qui appartient \`a la famille exponentielle et que les param\`etres de cette distribution sont $\theta_i$ et $\phi$. Encore une fois, On suppose que chaque observation poss\`de son propre param\`etre canonique $\theta_i$ mais que le param\`etre de dispersion $\phi$ est commun \`a l'ensemble des variables. Ainsi, on a $E[Y_i]=\mu_i=b'(\theta_i)$ et $Var[Y_i]=\phi b''(\theta_i) = \phi V(\mu_i)$.
\item[$\bullet$] {\bf Composante syst\'ematique:} On sp\'ecifie une composante syst\'ematique $\eta_i=f_\beta(X_{i1},\cdots,X_{ip})$  o\`u $f_\beta$ est une fonction connue param\'etr\'ee par un ensemble de param\`etres $\beta$. La composante syst\'ematique la plus populaire prend la forme lin\'eaire $\eta_i=\beta_0+\beta_1 X_{i1}+\beta_2 X_{i2} + \cdots + \beta_p X_{ip}$.
\item[$\bullet$] {\bf Fonction de lien:} La fonction de lien relie la composante al\'eatoire \`a la composante syst\'ematique $g(\mu_i)=\eta_i$. On en d\'eduit que $\eta_i=g[b'(\theta_i)]$. La fonction de lien $g(\cdot)$ est dite fonction de lien {\it canonique} si $g[b'(u)]=u, \forall u \in \Re$, autrement dit, si elle assure que $\eta_i=\theta_i$.
\end{enumerate}
\noindent Les param\`etres du mod\`ele consid\'er\'e sont $\beta$ et $\phi$.

\subsection{Estimation des param\`etres}

\noindent La fonction log-vraisemblanace des donn\'ees $\Delta=\{(Y_i,X_i), i=1,\cdots,n\}$ s'\'ecrit:
\begin{eqnarray*}
l(\beta,\phi|\Delta) &=& \sum_{i=1}^n \log\left[f_Y(y_i)\right] \\
&=& \sum_{i=1}^n \left\{ \log\left[a(y_i,\phi)\right] + \frac{y_i\theta_i-b(\theta_i)}{\phi} \right\}  \\
&=& \sum_{i=1}^n \log\left[a(y_i,\phi)\right] + \frac{1}{\phi} \sum_{i=1}^n \left[y_i\theta_i-b(\theta_i)\right] \\
&=& l_1(\phi) + \frac{1}{\phi} l_2(\beta),
\end{eqnarray*}
\noindent o\`u $l_1(\phi)=\sum_{i=1}^n \log\left[a(y_i,\phi)\right]$ et $l_2(\beta)=\sum_{i=1}^n \left[y_i\theta_i-b(\theta_i)\right]$

\noindent On calcule $\hat\beta$ l'estimateur de maximum de vraisemblance $\hat\beta$ en maximisant $l_2(\beta)$. On montre que la distribution asymptotique de cet estimateur est
$$\hat\beta \sim {\cal N}\left\{\beta,\phi(X^\top W X)^{-1}\right\},$$
\noindent o\`u $W=\mbox{diag}\{W_1,W_2,\cdots,W_n\}$ est une matrice diagonale de dimension $n$, $$W_i=\frac{1}{V(\hat\mu_i)\left[g'(\hat\mu_i)\right]^2},$$

\noindent et $\hat\mu_i=g^{-1}(\hat\eta_i)$ et $\hat\eta_i=\hat\beta_0+\hat\beta_1 X_{i1}+\cdots+\hat\beta_p X_{ip}$ sont les valeurs pr\'edites pour $\mu_i$ et $\eta_i$, respectivement.

\noindent Si $\phi$ est fix\'e, comme c'est le cas pour la r\'egression logistique ou la r\'egression Poisson, alors la proc\'edure d'estimation s'arr\^ete ici.


\noindent Si $\phi$ est inconnu, comme c'est le cas pour la r\'egression lin\'eaire avec une distribution normale ou Gamma, alors on doit l'estimer. L'estimateur de maximum de vraisemblance de $\phi$ est obtenu en maximisant $l_1(\phi) + l_2(\hat\beta)/\phi$. Cependant, l'estimateur obtenu est biais\'e pour les petites tailles d'\'echantillon. C'est pour cette raison qu'on lui pr\'ef\`ere des estimateurs obtenus par la m\'ethode des moments.

\noindent \`A cette fin, on calcule les statistiques 
$$X^2= \sum_{i=1}^n \frac{(y_i-\hat\mu_i)^2}{V(\hat\mu_i)}$$
et $$D(y,\hat \mu).$$

\noindent On montre que $X^2/\phi$ et $D(y,\hat\mu)/\phi$ sont approximativement distribu\'ees selon $\chi^2_{n-(p+1)}$. On en d\'eduit les estimateurs de Pearson et d\'eviance d\'efinis respectivement par $$\tilde \phi_P=\frac{X^2}{n-(p+1)}$$ et $$\tilde \phi_D=\frac{D(y,\hat\mu)}{n-(p+1)}.$$

\noindent Les deux estimateurs $\tilde \phi_P$ et $\tilde \phi_D$ sont faciles \`a calculer et poss\`edent de bonnes propri\'et\'es th\'eoriques. L'estimateur Pearson poss\`ede un biais n\'egligeable et est robuste lorsque la distribution est mal sp\'ecifi\'ee. C'est l'estimateur par d\'efaut de la fonction $\textsf{glm}$ du logiciel \textrm{R}.  L'estimateur d\'eviance poss\`ede une faible variance. Cependant, cet estimateur peut \^etre sensible lorsque certaines valeurs de $y$ sont proches des limites de leur support. C'est pour cette raison que $\tilde \phi_P$ est plus populaire.

\subsection{Inf\'erence statistique avec $\phi$ fix\'e}

\noindent Si $\phi$ est fix\'e, comme c'est le cas pour la r\'egression logistique ou la r\'egression Poisson, alors un intervalle de confiance \`a $(1-\alpha)100\%$ pour $\beta_j$ est construit comme suit:
$$\left[\hat\beta_j -z_{\alpha/2} \sqrt{ \phi v_j},\hat\beta_j +z_{\alpha/2} \sqrt{ \phi v_j} \right],$$
o\`u $v_j$ est $(j+1)^{\mbox{\`eme}}$ \'el\'ement de la diagonale de la matrice $(X^\top W X)^{-1}$, $j=0,1,\cdots,p$.

\noindent Consid\'erons maintenant une observation future $X^*=(1,X_1^*,X_2^*,\cdots,X_p^*)^\top$ et posons $\eta^*=\beta_0+\beta_1 X_1^* + \cdots + \beta_p X_p^*= {X^*}^\top\beta$ et $\mu^*=g^{-1}(\eta^*)$. La valeur pr\'edite de $\mu$ pour cette observation est $\hat\mu^*=g^{-1}(\hat\eta^*)$, o\`u $\hat\eta^*={X^*}^\top\hat\beta$. On  alors $Var[\hat\eta^*]={X^*}^\top Var[\hat\beta]X^* =\phi {X^*}^\top (X^\top W X)^{-1} X^*$. 

\noindent Des intervalles de confiance \`a $(1-\alpha)100\%$ pour $\eta^*$ et $\mu^*$ sont donc donn\'es par
$$\left[{X^*}^\top\hat\beta -  z_{\alpha/2} \sqrt{\phi {X^*}^\top (X^\top W X)^{-1} X^*}, {X^*}^\top\hat\beta +  z_{\alpha/2} \sqrt{\phi {X^*}^\top (X^\top W X)^{-1} X^*}\right]$$
et
$$\left[g^{-1}\left({X^*}^\top\hat\beta -  z_{\alpha/2} \sqrt{\phi {X^*}^\top (X^\top W X)^{-1} X^*}\right), g^{-1}\left({X^*}^\top\hat\beta +  z_{\alpha/2} \sqrt{\phi {X^*}^\top (X^\top W X)^{-1} X^*}\right)\right],$$
\noindent respectivement.

\noindent Supposons maintenant que l'on veuille comparer deux mod\`eles embo\^it\'es: un mod\`ele r\'eduit $A$ et un mod\`ele complet $B$. On construit la statistique de test
$$X_{obs} = \frac{D_A(y,\hat\mu)-D_B(y,\hat\mu)}{\phi}.$$

\noindent Cette statistique suit une loi de Khi-deux \`a $d$ degr\'es de libert\'e lorsque le mod\`ele r\'eduit est vrai. Ici, $d$ est la diff\'erence entre les nombres de param\`etres entre les deux mod\`eles. On rejette au seuil $\alpha$ l'hypoth\`ese que le mod\`ele r\'eduit est valide si $X_{obs}>\chi^2_{d,\alpha}$. Le seuil observ\'e est $\alpha^*=P(\chi^2_d > X_{obs})$.

\subsection{Inf\'erence statistique avec $\phi$ inconnu}

\noindent Lorsque le param\`etre $\phi$ est inconnu, on le remplace par son estimateur $\tilde \phi_P$. En faisant cette modification, on doit aussi remplacer les quantiles de la loi normale standard ${\cal N}(0,1)$ par les quantiles d'une loi $t$ de Student \`a $n-(p+1)$ degr\'es de libert\'e dans les expressions des bornes des intervalles de confiance de $\beta_j$, $\eta^*$ et $\mu^*$. Par exemple, l'intervalle de confiance pour $\beta_j$ devient
$$\left[\hat\beta_j -t_{\alpha/2,n-(p+1)} \sqrt{ \tilde \phi_P v_j},\hat\beta_j +t_{\alpha/2,n-(p+1)} \sqrt{ \tilde \phi_P v_j} \right],$$

\noindent De la m\^eme mani\`ere, l'expression de $X_{obs}$ pour comparer deux mod\`eles embo\^it\'es devient
$$X_{obs} = \frac{\{D_A(y,\hat\mu)-D_B(y,\hat\mu)\}/d}{\tilde \phi_P}.$$

\noindent La distribution de $X_{obs}$ lorsque le mod\`ele r\'eduit est valide devient $F_{d,n-(p+1)}$. On rejette au seuil $\alpha$ l'hypoth\`ese que le mod\`ele r\'eduit est valide si $X_{obs}>F_{d,n-(p+1),\alpha}$. Le seuil observ\'e devient $\alpha^*=P[F_{d,n-(p+1)} > X_{obs}]$.

\section{S\'election et validation de mod\`eles}

Les principes des proc\'edures pas \`a pas {\it backward}, {\it forward} et {\it stepwise} vues pour la r\'egression lin\'eaire et r\'egression logistique demeurent valident pour n'importe quel mod\`ele lin\'eaire g\'en\'eralis\'e. Ces mod\`eles se basent sur le $AIC$ et le $BIC$ d\'efinis par:
\begin{eqnarray*}
AIC &=& -2 l(\hat\beta,\tilde\phi_P|\Delta) + 2 (p+1) \\
BIC &=& -2 l(\hat\beta,\tilde\phi_P|\Delta) + \log(n) (p+1) \\
\end{eqnarray*}
\noindent Pareillement, les outils de validation de mod\`eles et de d\'etection d'observations aberrantes et influentes tels que les r\'esidus, les DFBETAS et les DFFITS vus dans les chapitres pr\'ec\'edents dans le cadre de r\'egression lin\'eaire et/ou r\'egression logistique restent valides pour l'ensemble des mod\`eles lin\'eaires g\'en\'eralis\'es.

\section{Exemple: r\'egression Poisson}

Consid\'erons les donn\'ees $\Delta=\{(Y_i,X_i),i=1\cdots,n\}$. La variable r\'eponse $Y_i$ d\'enote un d\'enombrement. Par exemple, $Y_i$ peut \^etre le nombre de nouveaux cas d'une maladie dans une ville durant une ann\'ee ou le nombre d'absences d'un \'el\`eve durant la session. Par cons\'equent, on suppose que $Y_i \in \mathbb{N}$. D'autre part, $X_i=(X_{i1},X_{i2},\cdots,X_{in})^\top$ est un ensemble de $p$ variables explicatives.

\noindent La r\'egression Poisson suppose que la variable r\'eponse $Y_i$ suit une loi Poisson avec un param\`etre $\mu_i$. Pour $y_i \in \mathbb{N}$, on a $$P(Y_i=y_i)=e^{-\mu_i} \frac{\mu_i^{y_i}}{y_i!}.$$
\noindent En plus, on a $E[Y_i]=Var[Y_i]=\mu_i$.


\noindent Comme on l'a vu plus haut, la famille exponentielle inclut la loi Poisson. Dans ce qui suit, on utilise la fonction de lien canonique $g(\mu_i)=\eta_i=\log(\mu_i)$. Avec cette fonction de lien, on a 
\begin{eqnarray*}
\mu_i &=& e^{\beta_0+\beta_1 X_{i1}+\cdots+\beta_p X_{ip}} \\
&=& e^{\beta_0} \times \left\{e^{\beta_1}\right\}^{X_{i1}} \times \cdots \times \left\{e^{\beta_p}\right\}^{X_{ip}}
\end{eqnarray*}
\noindent Ainsi, si $X_{ij}$ augmente de 1, $\mu_i$ est multipli\'e par $e^{\beta_j}$. Si $\beta_j>0$, alors $e^{\beta_j}>1$ et $\mu_i$ croit s'il est multipli\'e par $e^{\beta_j}$.

\noindent Les param\`etres de ce mod\`eles sont donc les coefficients de r\'egression $\beta=(\beta_0,\beta_1,\cdots,\beta_p)^\top$. L'estimateur de maximum de vraisemblence est obtenu en maximisant la fonction log-vraisemblance
$$l(\beta) = \sum_{i=1}^n \left\{y_i(\beta_0+\beta_1 X_{i1}+\cdots+\beta_p X_{ip})-e^{\beta_0+\beta_1 X_{i1}+\cdots+\beta_p X_{ip}}\right\} + K, $$
\noindent o\`u $K=-\sum_{i=1}^n \log(y_i!)$ est une constante qui ne d\'epent pas de $\beta$.

\noindent La distribution asymptotique de $\hat\beta$ v\'erifie
$$\hat\beta \sim {\cal N}\left\{\beta,(X^\top WX)^{-1}\right\},$$
\noindent o\`u $W$ est une matrice diagonale de taille $n$ avec les \'el\'ements $w_i=\hat\mu_i=\sout{\hat\beta_0+\hat\beta_1 X_{i1}+\cdots+\hat\beta_p X_{ip}}\textcolor{red}{e^{\hat\beta_0+\hat\beta_1 X_{i1}+\cdots+\hat\beta_p X_{ip}}}, i=1,\cdots,n$ sur la diagonale.

\noindent La d\'eviance de ce mod\`ele est 
$$d(y,\mu)=2\left\{y\log(\frac{y}{\mu})-(y-\mu)\right\}$$
et on v\'erifie l'ad\'equation de ce mod\`ele \`a un ensemble de donn\'ees a l'aide des statistiques d'ajustement
$$X^2 = \sum_{i=1}^n \frac{(y_i-\hat\mu_i)^2}{\hat\mu_i} \, \, \, \mbox{ et } \, \, \, D = \sum_{i=1}^n d(y_i,\hat\mu_i).$$
\noindent Ces deux statistiques suivent approximativement une loi de Khi-deux \`a $n-(p+1)$ degr\'es de libert\'e si le mod\`ele de Poisson est bien adapt\'e aux donn\'ees. 

\noindent Les proc\'edures d'inf\'erence statistique et de s\'election et validation de mod\`eles sont conduites en suivant la m\'ethodologie g\'en\'erale pour les mod\`eles lin\'eaires g\'en\'eralis\'es pr\'esent\'ee plus haut.

\noindent Dans ce qui suit, on s'int\'eresse \`a deux aspects sp\'ecifiques 'a la r\'egression Poisson: les variables explicatives {\it offset} et l'extra-variabilit\'e, aussi appel\'ee variabilit\'e extra-Poissonni\`ene. 

\noindent Pour ces illustrer ces aspects, on consid\`ere le jeu de donn\'ees \textsf{danishlc} de la library \texttt{GLMsData} du logiciel \textrm{R}. Ce jeu de donn\'ees contient deux variables explicatives : la ville et la tranche d'\^age. En effet, les donn\'ees sont collect\'ees dans 4 villes diff\'erentes et concernent 6 tranches d'\^age. La variable r\'eponse est le nombre de nouveaux cas de cancer des poumons entre 1968 et 1971. La base de donn\'ees contient une autre variable int\'eressante: la population de chaque tranche d'\^age dans chaque ville. 

\subsection{Variables explicatives {\it offset}}

La variable r\'eponse $Y_i$ est le nombre de nouveaux cas de cancer des poumons parmi $m_i$ qui repr\'esente la population d'une tranche d'\^age dans une ville donn\'ee. Normalement, on aurait eu une r\'egression logistique avec $Y_i \sim \mbox{Bin}(m_i,\pi_i)$, o\`u $\pi_i$ est la probabilit\'e qu'un individu d\'eveloppe le cancer en question. Une estimation empirique de $\pi_i$ est donn\'ee par $\tilde\pi_i=Y_i/m_i$. Selon les donn\'ees qu'on a, $\sout{\tilde pi_1}\textcolor{red}{\tilde\pi_1}=11/3059 \approx 3.6 \times 10^{-3}$. Cette est probabilit\'e est tr\`es petite et une r\'egression logistique ne semble pas \^etre appropri\'ee pour analyser ces donn\'ees. 

\noindent Une alternative est de mod\'eliser l'esp\'erance du ratio $Y_i/m_i$ en fonction des variables explicatives. Autrement dit, on \'ecrit que $Y_i$ suit une loi de Poisson avec un param\`etre $\mu_i$ qui v\'erifie
$$\log(\frac{\mu_i}{m_i}) = \beta_0+\beta_1 X_{i1}+\cdots+\beta_p X_{ip}.$$
\noindent Cette derni\`ere \'egalit\'e s'\'ecrit aussi:
$$\log(\mu_i) = \beta_0+\log(m_i)+\beta_1 X_{i1}+\cdots+\beta_p X_{ip}.$$
\noindent Ainsi, le terme $\log(m_i)$ peut \^etre consid\'er\'e comme une variable explicative avec un coefficient $\beta_j$ fix\'e \`a 1. On parle d'une variable explicative {\it offset}. Dans une r\'egression Poisson avec un lien $\log$, un terme {\it offset} est utile lorsqu'on pense que l'esp\'erance de la variable $Y_i$ est proportionnelle \`a une autre variable qu'on veut inclure dans le mod\`ele.
 
\begin{verbatim}
> if (!require("GLMsData")) install.packages("GLMsData")
> library("GLMsData")#### Installer la librairie GLMsData
> 
> data(danishlc)  #### Rendre le jeu de donnees danishlc disponible
> 
> modele.additif<-glm(Cases~offset(log(Pop))+City+Age,family=poisson,data=danishlc)
> modele.City<-glm(Cases~offset(log(Pop))+City,family=poisson,data=danishlc)
> modele.Age<-glm(Cases~offset(log(Pop))+Age,family=poisson,data=danishlc)
> 
> Q1 <- deviance(modele.City)-deviance(modele.additif)
> Q2 <- deviance(modele.Age)-deviance(modele.additif)
> 
> df1 <- summary(modele.City)$df[2]-summary(modele.additif)$df[2]
> df2 <- summary(modele.Age)$df[2]-summary(modele.additif)$df[2]
> 
> 1-pchisq(Q1,df=df1)
[1] 0
> 1-pchisq(Q2,df=df2)
[1] 0.182414
\end{verbatim} 

\noindent Il y a une diff\'erence significative entre les incidences du cancer des poumons des diff\'erentes tranches d'\^age. Cependant, ces incidences par tranche d'\^age 
 semblent \^etre les m\^emes pour les diff\'erentes villes.
 
\subsection{Variabilit\'e extra-Poissonni\`ene}

Avec une variable r\'eponse $Y_i$ qui suit une distribution Poisson, on doit avoir $E[Y_i]=Var[Y_i]=\mu_i$. Dans la pratique, il arrive qu'on ait $\sout{E[Y_i]>Var[Y_i]}\textcolor{red}{Var[Y_i]>E[Y_i]}$. Dans ce qui suit, on pr\'esente deux strat\'egies pour rem\'edier \`a ce probl\`eme. 

\subsubsection{Approche Quasi-Poisson} La premi\`ere strat\'egie consiste \`a consid\'erer un mod\`ele Quasi-Poisson. Avec ce mod\`ele, on sp\'ecifie seulement les deux premiers moment de la distribution de $Y_i$: $E[Y_i]=\mu_i$ et $Var[Y_i]=\phi \mu_i$. Avec cette approche, on obtient les m\^emes estimateurs $\hat\beta$ qu'un mod\`ele de Poisson mais la matrice $Var(\hat\beta)$ est multipli\'e par $\tilde \phi_P=X^2/(n-(p+1))$.

\begin{verbatim}
> ############################
> ## Approche Quasi-Poisson ##
> ############################
> 
> modele.Age.Quasi <- glm(Cases~offset(log(Pop))+Age,
+                         family=quasipoisson,data=danishlc)
> Pearson <- sum((danishlc$Cases-fitted(modele.Age))^2/fitted(modele.Age))
> phi.P.tilde <- Pearson/(n-(p+1)); phi.P.tilde
[1] 1.447135
> summary(modele.Age.Quasi)

Call:
glm(formula = Cases ~ offset(log(Pop)) + Age, family = quasipoisson, 
    data = danishlc)

Deviance Residuals: 
    Min       1Q   Median       3Q      Max  
-2.8520  -0.6424  -0.1067   0.7853   1.5468  

Coefficients:
            Estimate Std. Error t value Pr(>|t|)    
(Intercept) -4.45397    0.21606 -20.615 5.71e-14 ***
Age40-54    -1.40828    0.30089  -4.680 0.000186 ***
Age55-59    -0.32594    0.30316  -1.075 0.296511    
Age60-64     0.09339    0.28344   0.330 0.745574    
Age65-69     0.34201    0.28079   1.218 0.238937    
Age70-74     0.43894    0.28785   1.525 0.144669    
---
Signif. codes:  0 ‘***’ 0.001 ‘**’ 0.01 ‘*’ 0.05 ‘.’ 0.1 ‘ ’ 1

(Dispersion parameter for quasipoisson family taken to be 1.447135)

    Null deviance: 129.908  on 23  degrees of freedom
Residual deviance:  28.307  on 18  degrees of freedom
AIC: NA

Number of Fisher Scoring iterations: 5

> cbind(coef(modele.Age),coef(modele.Age.Quasi))
                   [,1]        [,2]
(Intercept) -4.45397213 -4.45397213
Age40-54    -1.40828068 -1.40828068
Age55-59    -0.32593887 -0.32593887
Age60-64     0.09339483  0.09339483
Age65-69     0.34200597  0.34200597
Age70-74     0.43894138  0.43894138
> phi.P.tilde*vcov(modele.Age)
            (Intercept)    Age40-54    Age55-59    Age60-64    Age65-69    Age70-74
(Intercept)  0.04668176 -0.04668176 -0.04668176 -0.04668176 -0.04668176 -0.04668176
Age40-54    -0.04668176  0.09053432  0.04668176  0.04668176  0.04668176  0.04668176
Age55-59    -0.04668176  0.04668176  0.09190471  0.04668176  0.04668176  0.04668176
Age60-64    -0.04668176  0.04668176  0.04668176  0.08033605  0.04668176  0.04668176
Age65-69    -0.04668176  0.04668176  0.04668176  0.04668176  0.07884030  0.04668176
Age70-74    -0.04668176  0.04668176  0.04668176  0.04668176  0.04668176  0.08286012
> vcov(modele.Age.Quasi)
            (Intercept)    Age40-54    Age55-59    Age60-64    Age65-69    Age70-74
(Intercept)  0.04668176 -0.04668176 -0.04668176 -0.04668176 -0.04668176 -0.04668176
Age40-54    -0.04668176  0.09053432  0.04668176  0.04668176  0.04668176  0.04668176
Age55-59    -0.04668176  0.04668176  0.09190471  0.04668176  0.04668176  0.04668176
Age60-64    -0.04668176  0.04668176  0.04668176  0.08033605  0.04668176  0.04668176
Age65-69    -0.04668176  0.04668176  0.04668176  0.04668176  0.07884030  0.04668176
Age70-74    -0.04668176  0.04668176  0.04668176  0.04668176  0.04668176  0.08286012
\end{verbatim}

\noindent Cette approache permet de calculer $\hat\beta$ et $Var[\hat\beta]$, de contruire des intervalles de confiance et d'effectuer des tests d'hypoth\`eses sur les coefficients $\beta$. Cependant, il est primordial de pr\'eciser que cette approche n'implique pas de distribution et donc ne permet pas de calculer des quantit\'es comme $AIC$ ou $BIC$ car on ne d\'efinit pas de fonction de log-vraisemblance.

\subsubsection{Mod\`ele avec la loi binomiale n\'egative}

On dit que $Y_i$ suit une loi binomiale n\'egative avec les param\`etres $\mu_i$ et $\alpha$ si 
$$P(Y_i=y_i) = \frac{\Gamma(y_i+1/\alpha)}{\Gamma(y_i+1)\Gamma(1/\alpha)} \left(\frac{1/\alpha}{\mu_i+1/\alpha}\right)^{1/\alpha} \left(\frac{\mu_i}{\mu_i+1/\alpha}\right)^{y_i}, \, \, \, y_i=0,1,2,\cdots, \, \, \, \mbox{ et } \alpha>0.$$

\noindent Avec cette loi, on a $E[Y_i]=\mu_i$ et $Var[Y_i]=\mu_i+\alpha \mu_i^2$. On retrouve la loi de Poisson de param\`etre $\mu_i$ lorsque $\alpha \rightarrow 0$.

\noindent Pour d\'efinir un mod\`ele de r\'egression, on pose $\mu_i = e^{\beta_0+\beta_1 X_{i1}+\cdots+\beta_p X_{ip}}$. Avec cette fonction de lien, les coefficients $\beta$ gardent la m\^eme interpr\'etation que dans une r\'egression Poisson.

\noindent Les param\`etres de ce mod\`ele\sout{s} sont donc $\beta$ et $\alpha$ et on retrouve la r\'egress\sout{s}ion Poisson lorsque $\alpha \rightarrow 0$. Le principal inconv\'eniant de ce mod\`ele est qu'on doit estimer $\alpha$ et $\beta$ simultan\'ement. Ceci peut causer des probl\`emes au niveau num\'erique.

\noindent Un avantage int\'eressant pour ce mod\`ele est qu'il nous permet de tester l'hypoth\`ese nulle $H_0:\alpha=0$. \`A cet effet, on ajuste le mod\`ele Poisson et le mod\`ele avec une loi binomiale n\'egative et on calcule les d\'eviances des deux mod\`eles, disons $D_0$ et $D_1$. La statistique du test sera alors $Q_{obs}=D_0-D_1$. Sous l'hypoth\`ese nulle $H_0:\alpha=0$, on a $$Q \sim \frac{1}{2} \chi^2_0 + \frac{1}{2} \chi^2_1$$
car on teste $\alpha$ sur la limite de son domaine de d\'efinition. Le seuil observ\'e est alors \'egal \`a $\alpha^*=P(\chi^2_1>Q_{obs})/2$. 

\begin{verbatim}
> #################################
> ## Approche binomiale negative ##
> #################################
> 
> library(MASS)
> 
> modele.Age.NB <- glm.nb(formula = Cases~offset(log(Pop))+Age,link="log",data=danishlc)
> 
> summary(modele.Age.NB)

Call:
glm.nb(formula = Cases ~ offset(log(Pop)) + Age, data = danishlc, 
    link = "log", init.theta = 195.5841968)

Deviance Residuals: 
    Min       1Q   Median       3Q      Max  
-2.8126  -0.6293  -0.1091   0.7648   1.5098  

Coefficients:
            Estimate Std. Error z value Pr(>|z|)    
(Intercept) -4.45209    0.18300 -24.328  < 2e-16 ***
Age40-54    -1.41012    0.25510  -5.528 3.25e-08 ***
Age55-59    -0.32639    0.25689  -1.271   0.2039    
Age60-64     0.09184    0.24089   0.381   0.7030    
Age65-69     0.34153    0.23871   1.431   0.1525    
Age70-74     0.43705    0.24449   1.788   0.0738 .  
---
Signif. codes:  0 ‘***’ 0.001 ‘**’ 0.01 ‘*’ 0.05 ‘.’ 0.1 ‘ ’ 1

(Dispersion parameter for Negative Binomial(195.5842) family taken to be 1)

    Null deviance: 122.207  on 23  degrees of freedom
Residual deviance:  27.168  on 18  degrees of freedom
AIC: 138.67

Number of Fisher Scoring iterations: 1


              Theta:  196 
          Std. Err.:  1364 

 2 x log-likelihood:  -124.673 
> 
> Q.obs <- deviance(modele.Age)-deviance(modele.Age.NB)
> 
> pvalue <- (1-pchisq(Q.obs,df=1))/2 
> 
> pvalue
[1] 0.1429792
\end{verbatim}

\noindent On ne rejette pas l'hypoth\`ese nulle $H_0:\alpha=0$. On peut donc supposer qu'il n'y a pas de variabilit\'e extra-Poissonni\`ene dans les donn\'ees.

\end{document}
